% ==============================================================================
% T0 Theory: Document 161 (English)
% Title: T0 Theory and Coherent Ising Machines
% Author: Johann Pascher
% Date: February 2026
% ==============================================================================

\section*{Abstract}

We analyze the structural parallels between the T0 time-mass duality theory and the
Coherent Ising Machine (CIM), a photonic analog computer for solving combinatorial
NP-hard optimization problems. We show that both systems rest on the same fundamental
principle: \emph{optimization as physical relaxation into a geometric minimum} -- not
as algorithmic iteration. Six structural parallels are identified, covering determinism,
ground state, qubit structure, coupling matrix, fractal corrections, and toroidal
topology. The central consequence is the concept of a \emph{T0 Ising Machine}, in which
all couplings $J_{ij}$ are not externally programmed but geometrically emerge from the
intrinsic time field $\Tfield$.

\clearpage

\section{Introduction}
\label{sec:intro}

\subsection{Motivation}

The Coherent Ising Machine (CIM) is an optical analog computer that solves combinatorial
optimization problems by physical relaxation into the ground state of an Ising Hamiltonian
\cite{inagaki_cim_2016, mcmahon_cim_2016}. The fundamental principle -- the system
``falls'' physically into its solution -- stands in direct conceptual relation to the
T0 theory \cite{pascher_t0_2025}, which derives all fundamental constants from a single
geometric parameter $\xipar = 4/30000$.

This document investigates this connection systematically. We show that T0 not only
provides a conceptual parallel to the CIM, but enables a fundamentally stronger
formulation: a T0-based Ising Machine in which the coupling matrix $J_{ij}$ geometrically
emerges from the mass field $\mfield$ -- rather than being externally programmed.

\subsection{Structure}

Section~\ref{sec:ising} describes the Ising model as a universal optimization language.
Section~\ref{sec:cim} explains the operating principle of the CIM.
Section~\ref{sec:t0} summarizes the relevant T0 foundations.
Section~\ref{sec:parallels} systematically analyzes the six structural parallels.
Section~\ref{sec:t0machine} formulates the concept of the T0 Ising Machine.
Section~\ref{sec:questions} discusses open research questions.

\clearpage

\section{The Ising Model as Universal Optimization Language}
\label{sec:ising}

\subsection{The Hamiltonian}

The Ising model from solid-state physics describes magnetic spins
$\sigma_i \in \{-1, +1\}$ interacting through a Hamiltonian:

\begin{equation}
    H_{\text{Ising}} = -\frac{1}{2} \sum_{i \neq j} J_{ij} \, \sigma_i \sigma_j
    - \sum_{i} h_i \, \sigma_i
    \label{eq:ising_hamiltonian}
\end{equation}

Here $J_{ij}$ are the coupling strengths between spins $i$ and $j$, and $h_i$ are
external fields (bias terms). The ground state -- the spin configuration with minimal
$H_{\text{Ising}}$ -- corresponds to the optimal solution of the encoded problem.

\subsection{QUBO: The Universal Mapping}

Every NP-hard combinatorial problem can be formulated as minimization of the Ising
Hamiltonian \eqref{eq:ising_hamiltonian} via \emph{Quadratic Unconstrained Binary
Optimization} (QUBO):

\begin{table}[h]
    \centering
    \caption{Examples of the mapping of NP-hard problems onto Ising Hamiltonians}
    \label{tab:qubo_examples}
    \resizebox{\textwidth}{!}{\begin{tabular}{@{}ll@{}}
        \toprule
        Optimization Problem & Ising Correspondence \\
        \midrule
        Travelling Salesman Problem & Couplings $J_{ij}$ as distance matrix \\
        MaxCut (graph partitioning) & Spins represent node assignments \\
        Portfolio optimization & Spins as stock selection (buy/not buy) \\
        Protein structure prediction & Spins as torsion angle configurations \\
        RSA factorization & Spins as bit representation of factors \\
        \bottomrule
    \end{tabular}}
\end{table}

This mapping is not an approximation -- it is exact. Finding the ground state of
\eqref{eq:ising_hamiltonian} solves the optimization problem.

\begin{important}[Central Principle]
    Combinatorial optimization $\equiv$ search for the ground state of the Ising
    Hamiltonian. This is the conceptual core of both the CIM and -- as we will show --
    the T0 theory.
\end{important}

\clearpage

\section{Coherent Ising Machine: Operating Principle}
\label{sec:cim}

\subsection{Physical Implementation}

A Coherent Ising Machine (CIM) physically implements the Ising model through a
network of time-multiplexed degenerate optical parametric oscillators (DOPO)
\cite{yamamoto_cim_2017, honjo_100k_cim_2021}:

\begin{itemize}
    \item \textbf{Spin} $\sigma_i = +1$: laser pulse with phase $0°$
    \item \textbf{Spin} $\sigma_i = -1$: laser pulse with phase $180°$
    \item \textbf{Coupling} $J_{ij}$: electronic feedback network between pulses
    \item \textbf{Dynamics}: physical relaxation into the energy minimum
\end{itemize}

The system evolves deterministically according to the field equations of the
parametric oscillator -- and in doing so automatically finds (or with high probability
finds) the ground state of the Hamiltonian \eqref{eq:ising_hamiltonian}.

\subsection{Scaling and State of the Art}

\begin{table}[h]
    \centering
    \caption{Comparison: Classical computer vs.\ Coherent Ising Machine}
    \label{tab:cim_vs_classical}
    \resizebox{\textwidth}{!}{\begin{tabular}{@{}lll@{}}
        \toprule
        Property & Classical Computer & Coherent Ising Machine \\
        \midrule
        Parallelism & Sequential / limited & Massively parallel (all spins simultaneous) \\
        Solution paradigm & Algorithm iterates & Physical relaxation \\
        Scaling & Exponentially harder & Natural via light propagation \\
        Energy demand & High for NP-hard problems & Potentially orders of magnitude lower \\
        Spins (as of 2024) & -- & Up to 100,000 optical spins \\
        \bottomrule
    \end{tabular}}
\end{table}

\subsection{Determinism Behind Apparent Stochasticity}
\label{subsec:cim_determinism}

A CIM appears to search stochastically: different runs sometimes yield different
solutions. But this is a surface property: physically the system follows
\emph{deterministically} the field equations of the OPO. The apparent randomness
is solely ignorance of the exact initial conditions of the quantum vacuum.

This is the first and most fundamental parallel to the T0 theory.

\clearpage

\section{T0 Foundations: Relevant Concepts}
\label{sec:t0}

\subsection{The Fundamental Binding Condition}

The T0 theory postulates time-mass duality as a fundamental geometric constraint:

\begin{equation}
    T \cdot m = 1 \quad \text{(in natural units)}
    \label{eq:t0_constraint}
\end{equation}

The intrinsic time field of a particle with mass $m$ is:

\begin{equation}
    \Tfield = \frac{\hbar}{m(x,t) \cdot c^2}
    \label{eq:t0_timefield}
\end{equation}

Time and mass are not independent quantities -- they are reciprocal aspects of a
single geometric reality.

\subsection{The Geometric Origin Parameter}

All fundamental constants emerge from a single dimensionless geometric parameter
\cite{pascher_xi_2025}:

\begin{equation}
    \xipar = \frac{4}{30000} \approx 1.333 \times 10^{-4}
    \label{eq:xi_param}
\end{equation}

This parameter defines the ratio between the Planck scale and the cosmic scale.
There are \textbf{no free parameters} -- the ground state of the theory is fully
determined by $\xipar$.

\subsection{T0 Qubits}

In the T0 theory, qubits emerge naturally from time-mass duality
\cite{pascher_quantum_2025}:

\begin{equation}
    |\text{T0-Qubit}\rangle = \alpha \, |T{\uparrow}, m{\downarrow}\rangle
    + \beta \, |T{\downarrow}, m{\uparrow}\rangle
    \label{eq:t0_qubit}
\end{equation}

with the binding condition $T \cdot m = 1$. The spin states are not chosen
arbitrarily -- they are forced by fundamental geometry:

\begin{itemize}
    \item $|T{\uparrow}, m{\downarrow}\rangle$: high time / low mass
    \item $|T{\downarrow}, m{\uparrow}\rangle$: low time / high mass
\end{itemize}

\subsection{The Fractal Correction}

The fractal correction $\Kfrak$ plays a central role in T0:
$\Kfrak^{3/2}$ explains the $\sim 2\%$ deviation in the anomalous magnetic
moment of the muon (g-2 anomaly) \cite{pascher_g2_2025}. This is structurally
analogous to the annealing mechanism in the CIM.

\clearpage

\section{Six Structural Parallels: T0 and CIM}
\label{sec:parallels}

\subsection{Parallel 1: Determinism Behind Apparent Stochasticity}

\begin{table}[h]
    \centering
    \caption{Parallel 1: Determinism}
    \label{tab:parallel1}
    \resizebox{\textwidth}{!}{\begin{tabular}{@{}lll@{}}
        \toprule
        Aspect & Coherent Ising Machine & T0 Theory \\
        \midrule
        Surface & Apparently stochastic search & Apparently probabilistic QM \\
        Deep structure & Deterministic (OPO field eqs.) & Deterministic ($\mfield$ geometry) \\
        Source of ``randomness'' & Unknown initial conditions & Unknown field configuration \\
        Consequence & Solution is physically forced & Measurement result geometrically fixed \\
        \bottomrule
    \end{tabular}}
\end{table}

Both systems replace fundamental probabilism with hidden determinism:
in the CIM the solution is not computed, it \emph{emerges}.
In T0 the measurement result is not chosen randomly -- it is already determined
by the mass field $\mfield$.

\begin{keyresult}[Parallel 1]
    Apparent probabilism $\equiv$ hidden determinism through geometric field structure --
    in CIM \emph{and} T0.
\end{keyresult}

\subsection{Parallel 2: Ground State as Geometric Minimum}

The CIM solves optimization through physical relaxation:
\begin{equation}
    H_{\text{Ising}} \rightarrow \text{minimum} \equiv \text{optimization solution}
    \label{eq:cim_groundstate}
\end{equation}

T0 describes the universe as already residing in the ground state of its own geometry:
\begin{equation}
    \xipar = \frac{4}{30000} \equiv \text{geometric minimum} \Rightarrow
    \text{all fundamental constants}
    \label{eq:t0_groundstate}
\end{equation}

In both cases the ``solution'' is not a computation -- it is the physically natural
resting state of the system. \textbf{Optimization is relaxation, not iteration.}

\subsection{Parallel 3: Qubit Structure}

\begin{table}[h]
    \centering
    \caption{Parallel 3: Comparison of qubit structures}
    \label{tab:parallel3}
    \resizebox{\textwidth}{!}{\begin{tabular}{@{}lll@{}}
        \toprule
        Aspect & CIM Spin & T0 Qubit \\
        \midrule
        Carrier & Phase of laser pulse & Time-mass relation \\
        Spin $+1$ & Phase $0°$ & $|T{\uparrow}, m{\downarrow}\rangle$ \\
        Spin $-1$ & Phase $180°$ & $|T{\downarrow}, m{\uparrow}\rangle$ \\
        Constraint & External (laser pump) & Intrinsic: $T \cdot m = 1$ \\
        Superposition & Coherent phase superposition & $\alpha|T{\uparrow},m{\downarrow}\rangle
        + \beta|T{\downarrow},m{\uparrow}\rangle$ \\
        Collapse & Phase measurement & Geometric field projection \\
        \bottomrule
    \end{tabular}}
\end{table}

\begin{important}[Fundamental Difference]
    The T0 spin is \emph{more fundamental} than the CIM spin: it carries its own
    constraint ($T \cdot m = 1$) intrinsically. A CIM spin must be externally
    constrained to $\pm 1$; a T0 spin is constrained by the geometry of spacetime itself.
\end{important}

\subsection{Parallel 4: The Coupling Matrix $J_{ij}$}
\label{subsec:coupling}

This is the deepest and physically most significant parallel:

\begin{table}[h]
    \centering
    \caption{Parallel 4: Coupling matrix -- the decisive difference}
    \label{tab:parallel4}
    \resizebox{\textwidth}{!}{\begin{tabular}{@{}lll@{}}
        \toprule
        Aspect & Classical CIM & T0 Ising Machine (concept) \\
        \midrule
        Couplings $J_{ij}$ & Externally programmed & Geometrically emergent \\
        Degrees of freedom & $J_{ij}$ freely selectable & Uniquely fixed by field geometry \\
        Physical meaning & Mathematical auxiliary & Real geometric interactions \\
        Universality & Arbitrary problems encodable & Physically meaningful problems \\
        Advantage & Maximum flexibility & No encoding errors possible \\
        \bottomrule
    \end{tabular}}
\end{table}

In a T0-based Ising Machine, $J_{ij}$ is \emph{not} a free parameter -- it emerges
directly from the intrinsic time field:

\begin{equation}
    J_{ij}^{(\text{T0})} = \int \Tfield_i \cdot \Tfield_j \cdot G(x_i, x_j) \, d^3x
    \label{eq:t0_coupling}
\end{equation}

where $G(x_i, x_j)$ is the geometric propagator of the T0 mass field.

\subsection{Parallel 5: Fractal Correction as Annealing Mechanism}

\begin{table}[h]
    \centering
    \caption{Parallel 5: Fractal correction and annealing}
    \label{tab:parallel5}
    \resizebox{\textwidth}{!}{\begin{tabular}{@{}lll@{}}
        \toprule
        Aspect & Simulated Annealing (CIM) & T0 Fractal Correction \\
        \midrule
        Parameter & Temperature $T_{\text{ann}}$ & Correction factor $\Kfrak$ \\
        Function & Prevents local minima & Explains systematic deviations \\
        Effect & Cooling $\rightarrow$ ground state & Scaling $\rightarrow$ precision \\
        Example & MaxCut optimization & g-2 anomaly: $\Kfrak^{3/2}$ explains $\sim 2\%$ \\
        Mechanism & Thermal fluctuations & Fractal self-similarity \\
        \bottomrule
    \end{tabular}}
\end{table}

The fractal correction $\Kfrak^{3/2}$ in T0 plays the analogous role of the
annealing temperature parameter: it prevents the system from becoming trapped in
a local minimum of the geometric energy landscape.

\subsection{Parallel 6: Toroidal Topology as Optimal CIM Geometry}

T0 describes the universe as a 4D torus crystal \cite{pascher_torus_2025}.
This topology has direct implications for a T0-based Ising Machine:

\begin{itemize}
    \item \textbf{Toroidal boundary conditions}: no boundary problem in optimization --
    all spins have equally many neighbors
    \item \textbf{Periodic structure}: natural implementation of translation-invariant
    couplings
    \item \textbf{4D instead of 3D}: richer coupling structure than classical 2D/3D
    CIM arrays
    \item \textbf{Crystal symmetry}: discrete symmetry groups $\rightarrow$ natural
    spin lattices without external structure specification
\end{itemize}

\begin{keyresult}[Parallel 6]
    A T0 Ising Machine would operate on the \emph{natural geometry of the universe} --
    not on an artificially constructed graph.
\end{keyresult}

\clearpage

\section{The T0 Ising Machine: Conceptual Formulation}
\label{sec:t0machine}

\subsection{Overall Structure}

Based on the six identified parallels, a T0-based Ising Machine can be conceptually
defined:

\begin{table}[h]
    \centering
    \caption{Overall comparison: Classical CIM vs.\ T0 Ising Machine}
    \label{tab:t0machine}
    \resizebox{\textwidth}{!}{\begin{tabular}{@{}lll@{}}
        \toprule
        Component & Classical CIM & T0 Ising Machine \\
        \midrule
        Spin carrier & DOPO laser pulses (phase) & Time field states $\Tfield$ \\
        Spin values & $\pm 1$ (phase $0°/180°$) & $|T{\uparrow},m{\downarrow}\rangle$,
        $|T{\downarrow},m{\uparrow}\rangle$ \\
        Couplings & Externally programmed & Geometrically emergent from $\mfield$ \\
        Topology & Arbitrary graph & 4D torus crystal \\
        Ground state & Solution sought & Geometric fixed point $\xipar$ \\
        Annealing & External (temperature) & Internal ($\Kfrak$ correction) \\
        Determinism & Hidden deterministic & Explicitly deterministic \\
        \bottomrule
    \end{tabular}}
\end{table}

\subsection{Formal Correspondence}

The complete mapping between T0 geometry and the Ising formalism reads:

\begin{align}
    \sigma_i &\leftrightarrow \operatorname{sign}\!\left(\Tfield_i - \bar{T}\right)
    &&\text{[spin from time field]}
    \label{eq:t0_spin_map}\\
    J_{ij} &\leftrightarrow \langle \Tfield_i \cdot \Tfield_j \rangle_{\xipar}
    &&\text{[coupling from $\xipar$ geometry]}
    \label{eq:t0_j_map}\\
    h_i &\leftrightarrow \frac{\partial \xipar}{\partial \Tfield_i}
    &&\text{[external field from $\xipar$ gradient]}
    \label{eq:t0_h_map}\\
    H_{\text{Ising}} &\leftrightarrow E_{\text{T0}} = \int |\nabla \Tfield|^2 \, d^4x
    &&\text{[energy functionals identical]}
    \label{eq:t0_energy_map}
\end{align}

The minimum of the T0 energy functional \eqref{eq:t0_energy_map} is identical to
the ground state of the associated Ising Hamiltonian -- with the crucial difference
that in T0 all quantities ($J_{ij}$, $h_i$, $H$) emerge from a single parameter
$\xipar$.

\subsection{The Central Consequence}

\begin{revolutionary}[Fundamental Implication of the T0 Ising Machine]
    If T0 is correct, then every physically realizable optimization problem is already
    encoded as a geometric structure in the mass field $\mfield$. A T0 Ising Machine
    would not optimize -- it would \emph{read} a physical reality. The solution is not
    the result of a search, but the unveiling of an already-existing geometric truth.
\end{revolutionary}

\clearpage

\section{Open Questions and Research Directions}
\label{sec:questions}

\subsection{Theoretical Openness}

\begin{itemize}
    \item What is the explicit geometric propagator $G(x_i, x_j)$ in the T0 mass field?
    What symmetries does the torus topology enforce?

    \item Is there a natural restriction on which NP-hard problems are representable
    through T0 geometry -- or are all problems in principle encodable?

    \item How does $\Kfrak$ behave as an annealing parameter for different problem
    classes? Are there universal scaling laws?

    \item What is the connection between the torus structure of the T0 universe and
    optimal CIM topologies (toroidal vs.\ Möbius-band couplings)?
\end{itemize}

\subsection{Experimental Directions}

\begin{itemize}
    \item \textbf{Coupling structure comparison}: compare the coupling matrix of
    existing CIMs with T0 field propagator predictions. Are there systematic
    deviations pointing to geometric constraints?

    \item \textbf{Toroidal CIM topologies}: do toroidal CIM architectures show
    performance advantages over open graphs that go beyond combinatorial explanations?

    \item \textbf{T0 qubit test}: do the T0 qubits defined in \eqref{eq:t0_qubit}
    show characteristic signatures in existing photonic quantum computing experiments?

    \item \textbf{Testable prediction}: in a CIM with $N$ spins, the T0 correction
    factor $\mathcal{D}(N) = \exp(-\xipar \ln N / \Dfrak)$ should produce a measurable
    $\ln(N)$-scaling residual that is qualitatively distinct from standard error models
    ($1/N$ or $e^{-N}$).
\end{itemize}

\subsection{Connection to Scramblon Physics}

This document extends \cite{pascher_scramblons_2026}, which analyzed the T0 connection
to scramblon physics. Combining both approaches yields a complete picture:

\begin{itemize}
    \item \textbf{Scramblon theory}: dynamic description -- how quantum information
    scrambles in a CIM-like system
    \item \textbf{T0 Ising Machine}: geometric description -- which fundamental
    couplings follow from the spacetime structure
    \item \textbf{Synthesis}: scramblons describe the \emph{dynamics} of information
    spreading; T0 describes the \emph{geometry} of the coupling structure
\end{itemize}

\clearpage

\section{Summary}
\label{sec:summary}

We have identified and systematically analyzed six structural parallels between
Coherent Ising Machines and the T0 theory:

\begin{table}[h]
    \centering
    \caption{Summary: Six parallels T0 $\leftrightarrow$ CIM}
    \label{tab:summary}
    \resizebox{\textwidth}{!}{\begin{tabular}{@{}clll@{}}
        \toprule
        \# & Parallel & CIM & T0 \\
        \midrule
        1 & Determinism & Hidden (OPO) & Explicit (geometry) \\
        2 & Ground state & Solution sought & Fixed point $\xipar$ \\
        3 & Qubit structure & Phase $0°/180°$ & $T \cdot m = 1$ states \\
        4 & Couplings $J_{ij}$ & Externally programmed & Geometrically emergent \\
        5 & Annealing & External temperature & $\Kfrak$ correction \\
        6 & Topology & Arbitrary graph & 4D torus \\
        \bottomrule
    \end{tabular}}
\end{table}

T0 does not merely provide a conceptual parallel to the CIM -- it offers a
fundamentally deeper justification: optimization is not a computation executed on an
artificially constructed system, but physical relaxation into a geometric truth already
inscribed in the mass field of the universe.

\begin{keyresult}[Central Result]
    The formal correspondence
    \begin{equation}
        H_{\text{Ising}} \leftrightarrow E_{\text{T0}} = \int |\nabla \Tfield|^2 \, d^4x
    \end{equation}
    shows: Ising optimization and T0 energy minimization are mathematically equivalent
    formulations of the same geometric principle -- with the decisive difference that
    in T0 \emph{all} couplings emerge parameter-free from $\xipar = 4/30000$.
\end{keyresult}

\begin{thebibliography}{99}

\bibitem{inagaki_cim_2016}
T.~Inagaki et al.,
\newblock A coherent Ising machine for 2000-node optimization problems,
\newblock \emph{Science} \textbf{354}, 603 (2016).

\bibitem{mcmahon_cim_2016}
P.~L.~McMahon et al.,
\newblock A fully programmable 100-spin coherent Ising machine with all-to-all connections,
\newblock \emph{Science} \textbf{354}, 614 (2016).

\bibitem{yamamoto_cim_2017}
Y.~Yamamoto et al.,
\newblock Coherent Ising machines -- optical neural networks operating at the quantum limit,
\newblock \emph{npj Quantum Inf.} \textbf{3}, 49 (2017).

\bibitem{honjo_100k_cim_2021}
T.~Honjo et al.,
\newblock 100,000-spin coherent Ising machine,
\newblock \emph{Sci.\ Adv.} \textbf{7}, eabh0952 (2021).

\bibitem{pascher_t0_2025}
J.~Pascher,
\newblock T0 Theory: Time-Mass Duality as a Fundamental Principle,
\newblock T0 Documentation Series, Document 003 (2025).

\bibitem{pascher_xi_2025}
J.~Pascher,
\newblock The Geometric Origin Parameter $\xi = 4/30000$,
\newblock T0 Documentation Series, Document 009 (2025).

\bibitem{pascher_quantum_2025}
J.~Pascher,
\newblock T0 Quantum Computing: Qubits and Logic Gates,
\newblock T0 Documentation Series, Document 147 (2025).

\bibitem{pascher_g2_2025}
J.~Pascher,
\newblock Fractal Corrections and the g-2 Anomaly,
\newblock T0 Documentation Series, Document 018 (2025).

\bibitem{pascher_torus_2025}
J.~Pascher,
\newblock The Universe as a 4D Torus Crystal: Geometric Cosmology,
\newblock T0 Documentation Series, Document 026 (2025).

\bibitem{pascher_scramblons_2026}
J.~Pascher,
\newblock T0 Projection onto Scramblon Physics,
\newblock T0 Documentation Series, Document 148 (2026).

\end{thebibliography}
